\documentclass[12pt,a4paper]{article}

\usepackage[T2A]{fontenc}
\usepackage[utf8]{inputenc}
\usepackage[russian]{babel}
\usepackage{amsmath,amssymb,amsthm,mathtools}
\usepackage{geometry}
\usepackage{microtype}
\usepackage{enumitem}
\usepackage{xcolor}
\usepackage{tikz}
\usepackage{pgfplots}
\pgfplotsset{compat=1.18}
\usepackage{hyperref}

\geometry{left=25mm,right=20mm,top=22mm,bottom=24mm}
\hypersetup{
  colorlinks=true,
  linkcolor=blue!45!black,
  urlcolor=blue!55!black
}
\setlist{nosep}
\setlength{\parindent}{1.25em}
\setlength{\parskip}{0.15em}

\newtheorem{theorem}{Теорема}
\newtheorem{proposition}{Утверждение}
\newtheorem{lemma}{Лемма}
\theoremstyle{definition}
\newtheorem{definition}{Определение}
\theoremstyle{remark}
\newtheorem*{remark}{Замечание}

\DeclareMathOperator{\Sp}{Sp}
\DeclareMathOperator{\diag}{diag}
\DeclareMathOperator{\cond}{cond}
\newcommand{\CC}{\mathbb{C}}
\newcommand{\RR}{\mathbb{R}}
\newcommand{\ZZ}{\mathbb{Z}}
\newcommand{\eps}{\varepsilon}
\newcommand{\norm}[1]{\left\lVert #1\right\rVert}

\DeclareMathOperator{\ind}{ind}
\newcommand{\Conv}{\mathcal C}

\newcommand{\TT}{\mathbb T}
\newcommand{\BB}{\mathcal B}
\title{Двумерные операторы свёртки\\Проекционные методы}
\author{Конспект лекции}
\date{14 ноября 2026 года}

\begin{document}
\maketitle

\section{Свёртка на двумерной решётке}

После одномерных операторов Винера--Хопфа перейдём к свёрткам
на группе
\[
  \ZZ^2=\ZZ\times\ZZ=\{k=(k_1,k_2):k_1,k_2\in\ZZ\}.
\]
Рассматриваем банахово пространство
\[
  X=\ell^1(\ZZ^2)
  =\left\{x=(x_k)_{k\in\ZZ^2}:\norm{x}_1
    :=\sum_{k\in\ZZ^2}|x_k|<\infty\right\}.
\]
Для абсолютно суммируемого семейства
\[
  \sum_{k\in\ZZ^2}a_k
  =\sum_{k_1=-\infty}^{\infty}\sum_{k_2=-\infty}^{\infty}a_{k_1,k_2};
\]
абсолютная сходимость позволяет переставлять порядок суммирования.
Пространства $\ell^1(\ZZ^2)$, $\ell^1(\ZZ)$ и $\ell^1(\mathbb N)$
изометрически изоморфны: достаточно выбрать биекцию множеств
индексов и перенумеровать координаты.

\begin{remark}
Это утверждение о банаховых пространствах. Произвольная
перенумерация не сохраняет сложение индексов и потому не превращает
двумерную свёртку в одномерную.
\end{remark}

\begin{definition}
Для $a,x\in\ell^1(\ZZ^2)$ их \emph{свёрткой} называется
\[
  (a*x)_k=\sum_{s\in\ZZ^2}a_{k-s}x_s.
\]
Оператор $A:X\to X$, $Ax=a*x$, называется \emph{оператором свёртки},
а последовательность $a$ --- его \emph{ядром}.
\end{definition}

Свойства свёртки те же, что в одномерном случае:
коммутативность, ассоциативность, билинейность и единица
$\delta_0$, где $(\delta_0)_0=1$, остальные координаты нулевые.
Развернём оценку, обеспечивающую корректность определения:
\begin{align*}
  \norm{a*x}_1
  &\leqslant\sum_{k\in\ZZ^2}\sum_{s\in\ZZ^2}|a_{k-s}|\,|x_s|\\
  &=\sum_{s\in\ZZ^2}|x_s|\sum_{r\in\ZZ^2}|a_r|
   =\norm{a}_1\norm{x}_1.
\end{align*}
Таким образом, $A$ ограничен. Поскольку $A\delta_0=a$ и
$\norm{\delta_0}_1=1$, получаем точное равенство
\[
  \boxed{\norm{A}=\norm{a}_1}.
\]
Эта оценка и пояснение об изоморфизме разворачивают краткие
замечания первой страницы рукописи.

\begin{remark}[Свёртка и корреляция в обработке изображений]
У Гонсалеса и Вудса эти операции различаются направлением смещения
маски~\cite[\S\,3.4.2, с.~188--192]{GW2012}. Для вещественных
данных в обозначениях лекции
\[
  (Ax)_k=\sum_{r\in\ZZ^2}a_rx_{k-r},\qquad
  (C_ax)_k=\sum_{r\in\ZZ^2}a_rx_{k+r}.
\]
Первая формула совпадает с определением свёртки после замены $r=k-s$;
вторая задаёт пространственную корреляцию в соглашении учебника.
Чтобы получить свёртку алгоритмом корреляции, надо заменить
$a_r$ на $a_{-r}$: в двумерной маске это поворот на $180^\circ$.
Для центрально-симметричной маски $a_r=a_{-r}$ результаты совпадают.

Простая проверка реализации --- подать единичный импульс:
$A\delta_0=a$, тогда как $(C_a\delta_0)_k=a_{-k}$.
Термин <<маска свёртки>> сам по себе не гарантирует, что программа
использует именно первый вариант. Для комплексных сигналов
соглашение о сопряжении в корреляции нужно оговаривать отдельно.
\end{remark}

\section{Символ на торе и два индекса}

Положим
\[
  \Gamma=\{z\in\CC:|z|=1\},\qquad
  \TT^2=\Gamma\times\Gamma,\qquad
  t^k=t_1^{k_1}t_2^{k_2}.
\]
Тор $\TT^2$ --- декартово произведение двух окружностей.
Его точки можно параметризовать как
$(e^{i\theta_1},e^{i\theta_2})$, где каждый угол берётся по модулю $2\pi$.

\begin{definition}
\emph{Символ} оператора свёртки $A$ с ядром $a$ --- функция
\[
  a(t)=\sum_{k\in\ZZ^2}a_kt^k
  =\sum_{k_1,k_2\in\ZZ}a_{k_1,k_2}t_1^{k_1}t_2^{k_2},
  \qquad t\in\TT^2.
\]
\end{definition}

Для ядра и символа, как и в рукописи, используется одна буква.
Поскольку $|t^k|=1$, ряд сходится абсолютно и равномерно
на торе. В частности, символ непрерывен и
$\max_{\TT^2}|a(t)|\leqslant\norm{a}_1$.
Алгебру таких функций обозначим $W(\TT^2)$; её норма равна
$\norm{a}_W=\sum_k|a_k|$.

\begin{figure}[ht]
\centering
\begin{tikzpicture}[scale=1.05,>=stealth]
  \draw[blue!60!black,thick] (0,0) circle (0.75);
  \fill (0.45,0.6) circle (0.035) node[above right] {$t_1$};
  \node at (1.3,0) {$\times$};
  \draw[blue!60!black,thick] (2.6,0) circle (0.75);
  \fill (3.2,-0.45) circle (0.035) node[below right] {$t_2$};
  \node at (4.0,0) {$\longleftrightarrow$};
  \draw[thick] (5,-0.75) rectangle (6.5,0.75);
  \draw[->,blue!60!black,thick] (5.35,-0.75)--(6.05,-0.75);
  \draw[->,blue!60!black,thick] (5.35,0.75)--(6.05,0.75);
  \draw[->,blue!60!black,thick] (5,-0.35)--(5,0.35);
  \draw[->,blue!60!black,thick] (6.5,-0.35)--(6.5,0.35);
  \node at (5.75,0) {$\TT^2$};
  \node[below] at (0,-0.9) {$\Gamma$};
  \node[below] at (2.6,-0.9) {$\Gamma$};
  \node at (7.2,0) {$\longleftrightarrow$};
  \begin{axis}[
    at={(9.6cm,0cm)},anchor=center,
    width=3.6cm,height=2.6cm,
    hide axis,axis equal image,scale only axis,
    view={35}{28},
    domain=0:360,y domain=0:360,
    samples=49,samples y=25,
    z buffer=sort,
    colormap={torusblue}{
      color(0cm)=(blue!65!black);
      color(1cm)=(blue!18!white)
    },
    point meta min=-0.75,point meta max=0.75
  ]
    % Стандартное вложение тора: R=2, r=0.75.
    % x --- угол вдоль большой окружности, y --- вдоль сечения трубки.
    \addplot3[surf,shader=interp,draw=none,
      point meta=z]
      ({(2+0.75*cos(y))*cos(x)},
       {(2+0.75*cos(y))*sin(x)},
       {0.75*sin(y)});
  \end{axis}
  \node[below] at (9.6,-1.25) {Тор $\TT^2$};
\end{tikzpicture}
\caption{Тор как произведение двух окружностей, квадрат со склеенными
противоположными сторонами и его пространственная модель.
Один угол задаёт обход вокруг отверстия, другой --- вокруг сечения трубки.}
\end{figure}

\subsection*{Как символ на торе связан с двумерным ДПФ}

В учебнике прямое ДПФ массива $h$ размера $P\times Q$ записано
с отрицательным знаком в экспоненте, а множитель $1/(PQ)$ стоит
в обратном преобразовании~\cite[\S\,4.5.5, с.~284]{GW2012}:
\[
  H(u,v)=\sum_{p=0}^{P-1}\sum_{q=0}^{Q-1}
    h_{p,q}e^{-2\pi i(up/P+vq/Q)}.
\]
Связь с лекционным символом получается подстановкой
\[
  H(u,v)=a\bigl(e^{-2\pi i u/P},e^{-2\pi i v/Q}\bigr),
  \qquad 0\leqslant u<P,\quad0\leqslant v<Q,
\]
если $h$ получен периодизацией ядра $a$ по модулю $(P,Q)$.
Для конечного ядра, размещённого в этом массиве без наложения,
это обычное дополнение нулями с учётом положения начала координат.
При усечении бесконечного ядра ДПФ даёт отсчёты символа
\emph{усечённого}, а не исходного ядра.

Таким образом, непрерывный тор $\TT^2$ --- область определения
символа свёртки на $\ZZ^2$, а ДПФ использует конечную сетку на торе.
Периодическое продолжение массива соответствует группе
$\ZZ_P\times\ZZ_Q$; его периодичность разобрана
в~\cite[\S\,4.6.3, с.~285--287]{GW2012}.
Это пояснение связывает две системы обозначений: сам переход
от символа лекции к отсчётам ДПФ здесь выписан дополнительно.

Предположим теперь, что $a(t)\ne0$ на всём торе.
Определим два целых числа
\[
  \kappa_1=\ind_{t_1}a(t_1,1),\qquad
  \kappa_2=\ind_{t_2}a(1,t_2),\qquad
  \kappa=(\kappa_1,\kappa_2)\in\ZZ^2.
\]
Здесь индекс --- число оборотов соответствующей ненулевой
замкнутой кривой вокруг начала координат при положительном
обходе окружности.

\begin{proposition}
Индекс по одной переменной не зависит от фиксированного значения другой:
\[
  \ind_{t_1}a(t_1,t_2^0)=\kappa_1,\qquad
  \ind_{t_2}a(t_1^0,t_2)=\kappa_2
  \quad(t_1^0,t_2^0\in\Gamma).
\]
\end{proposition}

\begin{proof}
Соединим $1$ и $t_2^0$ непрерывным путём $\gamma(s)$ на окружности,
$0\leqslant s\leqslant1$. Функция
\[
  H(t_1,s)=a(t_1,\gamma(s))
\]
задаёт гомотопию ненулевых петель между $a(t_1,1)$ и $a(t_1,t_2^0)$.
Индекс при такой гомотопии не меняется. Для второй координаты
доказательство аналогично.
\end{proof}
В рукописи первое равенство оставлено как упражнение;
доказательство и симметричное утверждение добавлены при оформлении.

\begin{theorem}
Двумерный оператор свёртки $A$ на $\ell^1(\ZZ^2)$ обратим
тогда и только тогда, когда его символ не обращается в нуль на торе:
\[
  \boxed{A\text{ обратим}\quad\Longleftrightarrow\quad
  a(t)\ne0\quad\text{для всех }t\in\TT^2}.
\]
\end{theorem}

\begin{proof}
Развернём указанную в рукописи ссылку на двумерную теорему Винера.
Эта теорема утверждает: если $a\in W(\TT^2)$ и $a$ нигде не равна
нулю, то $1/a\in W(\TT^2)$. Пусть $b$ --- последовательность
коэффициентов $1/a$. Тогда $a*b=b*a=\delta_0$, и свёртка с ядром
$b$ является обратным оператором к $A$.

Обратно, если $A$ обратим, существует $x=A^{-1}\delta_0\in\ell^1$.
Переход к символам в равенстве $a*x=\delta_0$ даёт
$a(t)x(t)=1$. Следовательно, $a$ не имеет нулей.
\end{proof}

Доказательство самой двумерной теоремы Винера в исходнике
не приводится и здесь не воспроизводится. Обратимость полной
свёртки не требует $\kappa=0$: например, символ $t_1$ задаёт
обратимый двусторонний сдвиг и имеет индексы $(1,0)$.
Последний пример добавлен для различения полной свёртки и её сжатий.

\subsection*{Обратимость и устойчивость восстановления изображения}

В модели Гонсалеса и Вудса наблюдаемое изображение задаётся
как свёртка исходного изображения с искажающим ядром плюс шум.
В частотной области это равенство имеет вид
\[
  G=HF+N,\qquad
  \widehat F_{\rm inv}=\frac{G}{H}=F+\frac{N}{H}
  \quad\text{при }H\ne0.
\]
Здесь $F,G,N$ --- Фурье-образы исходного изображения, наблюдения
и шума, а $H$ --- частотная характеристика искажения
\cite[\S\,5.1, с.~367--368; \S\,5.7, с.~411--414]{GW2012}.
Малые значения $|H|$ усиливают шум, даже если точных нулей нет.
Поэтому качественное утверждение об обратимости надо отличать
от количественной оценки нормы обратного оператора.

Для полной свёртки лекционная теорема обеспечивает
$A^{-1}x=b*x$, где символ $b$ равен $1/a$; тогда
$\norm{A^{-1}}_{\ell^1\to\ell^1}=\norm{b}_1$.
Эта величина может быть большой. Равенство
$\norm{A^{-1}}=1/\min|a|$ здесь не утверждается:
пространство лекции --- $\ell^1$, а не гильбертово пространство.

Учебник предлагает учитывать статистику сигнала и шума
винеровским фильтром~\cite[\S\,5.8, с.~414--416, формула (5.8-2)]{GW2012}:
\[
  \widehat F_{\rm W}
  =\frac{\overline H\,S_f}{|H|^2S_f+S_\eta}\,G.
\]
Здесь $S_f,S_\eta$ --- энергетические спектры сигнала и шума;
формула используется там, где знаменатель положителен.
В учебнике предполагаются некоррелированность изображения и шума,
нулевое среднее хотя бы у одного из них и линейность оценки.
При этих предпосылках минимизируется среднеквадратичная ошибка.
Если отношение $S_\eta/S_f$ заменено положительной константой
$\lambda$, получается практическая форма
\[
  \widehat F_\lambda=\frac{\overline H}{|H|^2+\lambda}\,G.
\]
Она не делит непосредственно на малое $H$, но вносит смещение
в оценку сигнала. Это восстановление с учётом шума, а не точное
обращение произвольного оператора. Винеровский фильтр главы~5
и теорема Винера об обратимости в алгебре $W$ --- разные результаты.

\section{Координатные проекторы и сжатия}

Для $U\subseteq\ZZ^2$ положим
\[
  (P_Ux)_k=\begin{cases}x_k,&k\in U,\\0,&k\notin U.\end{cases}
\]
Это ограниченный проектор на замкнутое подпространство
\[
  X_U=\operatorname{Im}P_U
  =\{x\in\ell^1(\ZZ^2):x_k=0\text{ при }k\notin U\}.
\]
Непосредственно из координатного определения следуют
\begin{align*}
  &P_U^2=P_U,\qquad
   \norm{P_U}=\begin{cases}1,&U\ne\varnothing,\\0,&U=\varnothing;\end{cases}\\
  &P_UP_V=P_{U\cap V}=P_VP_U;\\
  &P_{U\cup V}=P_U+P_V-P_{U\cap V};\\
  &U\cap V=\varnothing\quad\Longrightarrow\quad
    P_U+P_V=P_{U\cup V};\\
  &P_U+P_{\ZZ^2\setminus U}=I.
\end{align*}
Например, формула для объединения --- это равенство индикаторов
$\mathbf1_{U\cup V}=\mathbf1_U+\mathbf1_V-\mathbf1_{U\cap V}$.

\begin{definition}
\emph{Сжатием} свёртки на $U$ называется оператор
\[
  T_U=(P_UAP_U)|_{X_U}:X_U\longrightarrow X_U.
\]
Если $U\subseteq\RR^2$ --- геометрическая область, под $P_U$
понимаем $P_{U\cap\ZZ^2}$.
\end{definition}

На всём $X$ оператор $P_UAP_U$ обнуляет $\ker P_U$.
Поэтому обратимость сжатия всегда означает обратимость на $X_U$.
Координатно
\[
  (T_Ux)_k=\sum_{s\in U}a_{k-s}x_s\quad(k\in U)
\]
для $U\subseteq\ZZ^2$. Это уточнение существенно во всех
последующих критериях.

\subsection*{Граница изображения: сжатие и периодическое продолжение}

Для конечного окна $U$ оператор $P_UAP_U$ означает:
сначала считать изображение нулевым вне $U$, затем выполнить
свёртку и оставить результат в $U$. Отражение у края, повтор
крайних пикселей и периодическое продолжение задают другие
операторы. Первые два варианта расширения также обсуждаются
в учебнике~\cite[\S\,3.4.2, с.~190]{GW2012}; их сопоставление
с $P_UAP_U$ --- пояснение в обозначениях лекции.

Для массивов $f,h$ на периодической сетке теорема о свёртке даёт
\[
  \operatorname{IDFT}\bigl(\operatorname{DFT}(f)
       \operatorname{DFT}(h)\bigr)=f*_{\rm circ}h.
\]
Здесь свёртка \emph{циклическая}: индексы берутся по модулю
размеров массива. При попытке вычислить таким способом обычную
свёртку без подготовки появляются наложения соседних периодов
\cite[\S\,4.6.6, с.~298--301]{GW2012}.

Если изображение имеет размеры $M\times N$, а конечное ядро ---
$m\times n$, для полной линейной свёртки достаточно дополнить
\emph{оба} массива нулями до общих размеров
\[
  \boxed{P\geqslant M+m-1,\qquad Q\geqslant N+n-1}.
\]
Это условия (4.6-29)--(4.6-30) учебника в наших обозначениях.
После обратного ДПФ выбирается нужная часть результата; для
сжатия --- окно $U$. Положение нулевого индекса ядра надо согласовать
с формулой свёртки, иначе результат дополнительно сдвинется.

\begin{remark}
Достаточное дополнение нулями устраняет \emph{наложение периодов}
при вычислении линейной свёртки. Оно не восстанавливает неизвестные
пиксели за границей кадра. Резкое обрезание изображения также
может вызывать спектральную утечку; сглаживающее окно уменьшает
резкость края, но изменяет исходные данные
\cite[\S\,4.6.6, с.~301, 304]{GW2012}.
Это разные эффекты, и их не следует смешивать.
\end{remark}

\section{Полуплоскости и углы}

Пусть $\ell$ --- прямая через начало координат, $\nu\ne0$ ---
её нормаль. Рассматриваем полуплоскость
\[
  \Pi=\{y\in\RR^2:\langle\nu,y\rangle\geqslant0\}
  \quad\text{или}\quad
  \Pi=\{y\in\RR^2:\langle\nu,y\rangle>0\}.
\]
То есть граница либо включена целиком, либо исключена.

\begin{theorem}[Симоненко; формулировка из лекции]
Пусть $a(t)\ne0$ для всех $t\in\TT^2$, а $\kappa$ --- вектор
индексов символа. Тогда
\[
  \boxed{T_\Pi\text{ обратим}\quad\Longleftrightarrow\quad
  \kappa\in\ell}
  \qquad\bigl(\text{то есть }\langle\nu,\kappa\rangle=0\bigr).
\]
\end{theorem}

Теорема в рукописи приведена без доказательства. Условие
ненулевого символа выписано явно: именно при нём определён
использованный вектор индексов. Следствия при этом условии:
\begin{enumerate}
  \item для верхней полуплоскости $T_\Pi$ обратим
        тогда и только тогда, когда $\kappa_2=0$;
  \item для правой полуплоскости $T_\Pi$ обратим
        тогда и только тогда, когда $\kappa_1=0$;
  \item сжатия обратимы для всех полуплоскостей с границей
        через начало тогда и только тогда, когда $\kappa=(0,0)$.
\end{enumerate}
Для последнего утверждения достаточно сначала взять верхнюю
и правую полуплоскости. Обратно, нулевой вектор лежит на любой
граничной прямой.

\begin{figure}[ht]
\centering
\begin{tikzpicture}[scale=0.9,>=stealth]
  \fill[blue!10] (-1.5,0) rectangle (1.5,1.6);
  \draw[->] (-1.7,0)--(1.8,0) node[right] {$k_1$};
  \draw[->] (0,-0.5)--(0,1.9) node[above] {$k_2$};
  \node at (0.8,1) {$\Pi$};
  \node[below] at (0,-0.65) {$\kappa_2=0$};
  \begin{scope}[xshift=4.8cm]
    \fill[blue!10] (0,0)--(2,0.45)--(0.7,2)--cycle;
    \draw[->] (-0.5,0)--(2.3,0) node[right] {$k_1$};
    \draw[->] (0,-0.5)--(0,2.3) node[above] {$k_2$};
    \draw[blue!60!black,thick] (2,0.45)--(0,0)--(0.7,2);
    \node at (0.85,0.95) {$K$};
    \node[below left] at (0,0) {$0$};
  \end{scope}
\end{tikzpicture}
\caption{Полуплоскость и угол с вершиной в начале координат}
\end{figure}

Теперь $K$ --- плоский выпуклый угол с вершиной в начале
координат и раствором между $0$ и $\pi$, как на рисунке лекции.
Рассматриваем $T_K=(P_KAP_K)|_{X_K}$.

\begin{theorem}[Необходимое условие из лекции]
Если $T_K$ обратим, то
\[
  a(t)\ne0\quad(t\in\TT^2),\qquad\kappa=(0,0).
\]
\end{theorem}

Это только необходимое условие. В рукописи отмечено, что
одномерная факторизационная схема не даёт здесь такого же
критерия обратимости. Следующая теорема характеризует другое
свойство --- нётеровость.

\begin{definition}
Пусть $Y$ --- банахово пространство. Оператор $B\in\BB(Y)$
называется \emph{нётеровым} (\emph{$\Phi$-оператором}), если
существует $R\in\BB(Y)$ такой, что
\[
  RB=I_Y+C_1,\qquad BR=I_Y+C_2,
\]
где $C_1,C_2$ --- компактные (вполне непрерывные) операторы.
Оператор $R$ называется регуляризатором.
\end{definition}

\begin{theorem}[Симоненко; критерий нётеровости из лекции]
Для рассматриваемого сжатия в угле
\[
  \boxed{T_K\text{ нётеров}\quad\Longleftrightarrow\quad
  a(t)\ne0\ (t\in\TT^2),\quad\kappa=(0,0)}.
\]
\end{theorem}

Теоремы об угле в исходнике даны без доказательств.
В определении для $T_K$ все операторы и единица действуют на $X_K$.

\begin{remark}
Обратимый оператор нётеров: можно взять $R=B^{-1}$ и $C_1=C_2=0$.
Однако обратимость по модулю компактных операторов не равносильна
обычной обратимости. В частности, из компактности $C$ не следует
обратимость $I+C$.

Для пояснения краткой записи рукописи рассмотрим ненулевой
конечномерный проектор $E$. Оператор $I-E$ нётеров, поскольку
при $R=I$ обе погрешности равны $-E$, но
$\ker(I-E)=\operatorname{Im}E\ne\{0\}$.
Это добавленный общий пример; он не утверждает, что $I-E$
является сжатием рассматриваемой свёртки.
\end{remark}

\section{Проекционные методы: сильная сходимость}

Далее $X$ --- произвольное банахово пространство, а
$P_n\in\BB(X)$ --- ограниченные проекторы: $P_n^2=P_n$.

\begin{definition}
Последовательность $P_n$ \emph{сильно сходится} к единице, если
\[
  P_n\xrightarrow{s}I
  \quad\Longleftrightarrow\quad
  \norm{P_nx-x}\longrightarrow0\quad\text{для каждого }x\in X.
\]
\end{definition}

Сходимость по операторной норме означает
$\norm{P_n-I}\to0$ и требует равномерной малости ошибки на
единичном шаре. Для нетождественных проекторов это невозможно.

\begin{proposition}
Если $P\ne I$ --- проектор, то $\norm{I-P}\geqslant1$.
Следовательно, $\norm{P_n-I}\to0$ для проекторов возможно
только тогда, когда $P_n=I$ при всех достаточно больших $n$.
\end{proposition}

\begin{proof}
$Q=I-P$ --- ненулевой проектор, поскольку $Q^2=I-2P+P^2=Q$.
Возьмём ненулевой $y\in\operatorname{Im}Q$. Тогда $Qy=y$, откуда
$\norm{Q}\geqslant\norm{Qy}/\norm{y}=1$.
\end{proof}

\begin{remark}[Дополнение: координатные усечения]
В $\ell^1(\mathbb N)$ рассмотрим усечения
\[
  P_nx=(x_1,\ldots,x_n,0,0,\ldots).
\]
Тогда
\[
  \norm{P_nx-x}_1=\sum_{j>n}|x_j|\longrightarrow0,
  \qquad \norm{I-P_n}=1.
\]
Таким образом, сильная сходимость есть, а сходимости по норме нет.
Аналогично усечения на квадраты $[-n,n]^2\cap\ZZ^2$
сильно сходятся к единице в $\ell^1(\ZZ^2)$.
\end{remark}

\begin{theorem}[Следствие теоремы Банаха--Штейнгауза]
Если $B_n\in\BB(X)$ и $B_n\xrightarrow{s}B$, то
\[
  \sup_n\norm{B_n}<\infty.
\]
В частности, из $P_n\xrightarrow{s}I$ следует
$C_P:=\sup_n\norm{P_n}<\infty$.
\end{theorem}

Действительно, для каждого $x$ последовательность $B_nx$
сходится и потому ограничена. Принцип равномерной ограниченности
превращает эту поточечную ограниченность в общую оценку норм.

\section{Приближённые уравнения и критерий применимости}

Пусть $A\in\BB(X)$ обратим и $P_n\xrightarrow{s}I$.
Обозначим $X_n=\operatorname{Im}P_n$ и
\[
  A_n=(P_nAP_n)|_{X_n}:X_n\to X_n.
\]
Вместо точного уравнения
\begin{equation}\label{eq:exact}
  Ax=b
\end{equation}
решаем проекционные уравнения
\begin{equation}\label{eq:approx}
  A_nx_n=P_nb,\qquad x_n\in X_n.
\end{equation}
Образы ограниченных проекторов замкнуты, поэтому $X_n$ ---
банаховы пространства с нормой, унаследованной от $X$.
Если $P_n$ конечного ранга, уравнение~\eqref{eq:approx}
является конечномерной системой. Конечный ранг для общего
критерия ниже не требуется.

\begin{definition}
Проекционный метод по системе $\{P_n\}$ \emph{применим} к $A$,
если существует $N$, такое что при всех $n>N$ и каждом $b\in X$
уравнение~\eqref{eq:approx} имеет единственное решение $x_n\in X_n$,
причём
\[
  \norm{x_n-A^{-1}b}\longrightarrow0.
\]
\end{definition}

\begin{theorem}[Критерий применимости проекционного метода]
Пусть $A$ обратим и $P_n\xrightarrow{s}I$. Метод применим
тогда и только тогда, когда существуют $N$ и $M<\infty$ такие, что
для всех $n>N$:
\begin{enumerate}
  \item $A_n:X_n\to X_n$ обратим;
  \item $\norm{A_n^{-1}}_{\BB(X_n)}\leqslant M$.
\end{enumerate}
\end{theorem}

\newpage
\subsection*{Доказательство критерия и оценка ошибки}

Рукопись заканчивается формулировкой этого критерия.
Следующее доказательство и оценка ошибки --- дополнение.

\begin{proof}
\emph{Достаточность.}
Положим $x=A^{-1}b$ и $x_n=A_n^{-1}P_nb$. Тогда
\begin{align*}
  A_n(x_n-P_nx)
  &=P_nb-P_nAP_nx\\
  &=P_nA(x-P_nx).
\end{align*}
Обе части лежат в $X_n$. Применяя обратный оператор, получаем
\[
  \norm{x_n-P_nx}
  \leqslant M\norm{P_n}\norm{A}\norm{(I-P_n)x}.
\]
По теореме Банаха--Штейнгауза $\norm{P_n}\leqslant C_P$, поэтому
\begin{equation}\label{eq:error}
  \boxed{\norm{x_n-x}
    \leqslant(1+MC_P\norm{A})\norm{(I-P_n)x}
    \longrightarrow0}.
\end{equation}
Единственность приближённого решения следует из обратимости $A_n$.

\emph{Необходимость.}
Пусть метод применим. При $n>N$ любая правая часть $y\in X_n$
представляется как $P_nb$ (можно взять $b=y$). Единственность
и существование решения означают биективность $A_n$.
По теореме об обратном операторе $A_n^{-1}$ ограничен.

Рассмотрим операторы решения на одном и том же пространстве:
\[
  S_n=A_n^{-1}P_n:X\to X_n\subseteq X,\qquad n>N.
\]
По применимости метода $S_nb\to A^{-1}b$ для каждого $b\in X$,
то есть $S_n\xrightarrow{s}A^{-1}$. Теорема Банаха--Штейнгауза
даёт $\sup_{n>N}\norm{S_n}<\infty$. Для $y\in X_n$
имеем $P_ny=y$, поэтому
\[
  \norm{A_n^{-1}y}=\norm{S_ny}\leqslant\norm{S_n}\norm{y}.
\]
Следовательно, $\norm{A_n^{-1}}\leqslant\norm{S_n}$,
и нормы обратных равномерно ограничены.
\end{proof}

\begin{remark}
Одна лишь обратимость каждого $A_n$ не заменяет равномерной оценки
норм $A_n^{-1}$. В критерии второе условие обеспечивает устойчивость
решения приближённых задач. Оценка~\eqref{eq:error} связывает
ошибку решения с ошибкой проекции точного решения.
\end{remark}

\subsection*{Связь с вычислительной устойчивостью}

В условиях критерия при возмущении правой части на $\delta b$
приближённое решение меняется на
\[
  \delta x_n=A_n^{-1}P_n\delta b,\qquad
  \norm{\delta x_n}\leqslant MC_P\norm{\delta b}.
\]
Эта оценка, выведенная из лекционного критерия, объясняет,
зачем нужен общий предел норм обратных: он ограничивает усиление
ошибки данных независимо от номера приближения. Сходную проблему
усиления шума при инверсной фильтрации показывает учебник
\cite[\S\,5.7, с.~412]{GW2012}.

\emph{Дополнительный пример для различения конечной сетки и тора.}
Возьмём символ $a(t_1,t_2)=1+t_1$. Он имеет нули при $t_1=-1$,
поэтому полная свёртка на $\ell^1(\ZZ^2)$ необратима.
Но на сетке ДПФ с нечётным числом $P$ узлов по первой координате
все значения
\[
  H(u,v)=1+e^{-2\pi i u/P}
\]
ненулевые, и соответствующая конечная \emph{циклическая} свёртка
$C_{P,Q}$ обратима. При этом
\[
  \min_{u,v}|H(u,v)|=2\sin\frac{\pi}{2P}\longrightarrow0,
  \qquad
  \norm{C_{P,Q}^{-1}}_2
  =\frac{1}{2\sin(\pi/(2P))}\longrightarrow\infty.
\]
Здесь индекс $2$ обозначает операторную норму для евклидовой
нормы конечных массивов. Равенство следует из унитарной
диагонализации циклической свёртки матрицей ДПФ; к норме
$\ell^1$ это равенство не переносится. При чётном $P$ сетка
содержит $t_1=-1$, и матрица уже вырождена.

Пример получен в конспекте из формул ДПФ и теоремы о свёртке
\cite[\S\,4.5.5, с.~284; \S\,4.6.6, с.~298]{GW2012},
а не переписан из учебника. Он показывает, почему проверки
ненулевых значений на одной конечной сетке недостаточно для
вывода об обратимости на всём торе. Циклические матрицы здесь
не отождествляются со сжатиями $A_n=P_nAP_n$: у них другие
граничные условия, поэтому это иллюстрация устойчивости,
а не применение проекционного критерия к периодическим матрицам.

\begin{thebibliography}{1}
\bibitem{GW2012}
Гонсалес Р., Вудс Р.
\emph{Цифровая обработка изображений.}
3-е изд., исправленное и дополненное.
Пер. с англ. Л.~И.~Рубанова, П.~А.~Чочиа;
под ред. П.~А.~Чочиа. М.: Техносфера, 2012. 1104~с.
ISBN 978-5-94836-331-8.
\end{thebibliography}

\noindent\emph{О ссылках.} Указаны печатные страницы этого издания.
В использованном PDF учебника номер страницы
просмотрщика на единицу больше печатного. Пояснения со ссылками
на учебник добавлены при оформлении и не являются расшифровкой
отдельных записей лекции. Теоремы Симоненко и проекционный критерий
сохранены как материал лекции; их доказательство учебнику
Гонсалеса и Вудса не приписывается.

\end{document}
